% Prostudováno: ucebnice/02_Funkce_a_rovnice/Linearni_funkce_a_rovnice/Linerání rovnice.tex,
% pododdíl Lineární rovnice s parametrem (ř. 2023–2151).
% Realisticky.cz: externi_zdroje/realisticky/02_Funkce_a_rovnice/08_Rovnice_a_nerovnice_s_parametrem/020801_Linearni_rovnice_s_parametrem_I/020801_Lekce.pdf
% a 020802_Linearni_rovnice_s_parametrem_II/020802_Lekce.pdf; využito rozdělení podle nulového koeficientu a kontroly zvláštní hodnoty. Úlohy vlastní.
\documentclass[11pt,a4paper]{article}
% Spolecny styl pro ucitelske pripravy MAT (Work/MAT)
% Verze 1.1, 2026-10-23
\NeedsTeXFormat{LaTeX2e}
\ProvidesPackage{priprava}[2026/10/23 v1.1 Sablona ucitelskych priprav MAT]

% Zakladni typografie a cestina
\RequirePackage[T1]{fontenc}
\RequirePackage[utf8]{inputenc}
\RequirePackage[czech]{babel}
\RequirePackage{lmodern}
\RequirePackage{microtype}
\RequirePackage{amsmath,amssymb}
\RequirePackage{xcolor}
\RequirePackage{enumitem}
\RequirePackage[most]{tcolorbox}
\RequirePackage{tabularx,booktabs,array}
\RequirePackage{fancyhdr}
\RequirePackage{needspace}
\RequirePackage{geometry}
\RequirePackage[hidelinks]{hyperref}

\geometry{a4paper,top=17mm,bottom=18mm,left=23mm,right=23mm}
\definecolor{notegray}{gray}{0.39}
\setlength{\parindent}{0pt}
\setlength{\parskip}{2.5pt plus 1pt minus 0.5pt}
\setlist[itemize]{topsep=2pt,itemsep=1pt,parsep=0pt,leftmargin=1.45em}
\setlist[enumerate]{topsep=2pt,itemsep=2pt,parsep=0pt,leftmargin=1.65em}
\setcounter{secnumdepth}{0}
\renewcommand{\arraystretch}{1.12}

% Zahlavi: vlevo rocnik | globalni tema, vpravo cislo: nazev hodiny
% Zapati zustava zcela prazdne.
\pagestyle{fancy}
\fancyhf{}
\renewcommand{\headrulewidth}{0.35pt}
\renewcommand{\footrulewidth}{0pt}
\setlength{\headheight}{14pt}
\newcommand{\nastavPripravu}[4]{%
  \fancyhead[L]{\small #1\enspace\textbar\enspace #2}%
  \fancyhead[R]{\small #3: #4}%
}

% Minimalisticke cernobile boxy. Kazda definice zvlast.
\tcbset{commonlesson/.style={enhanced,breakable=false,colback=white,colframe=black,
  colbacktitle=white,coltitle=black,boxrule=0.45pt,
  arc=0.8mm,outer arc=0.8mm,left=3mm,right=3mm,top=1.6mm,bottom=1.6mm,
  toptitle=1mm,bottomtitle=0.7mm,fonttitle=\bfseries,
  attach boxed title to top left={xshift=2.5mm,yshift=-1.15mm},
  boxed title style={colback=white,colframe=white,boxrule=0pt,left=1mm,right=1mm}}}
\newtcolorbox{zadaniBox}[1][]{commonlesson,title={Zadání},#1}
\newtcolorbox{prikladBox}[1][]{commonlesson,title={Příklad},#1}
\newtcolorbox{definiceBox}[1][]{commonlesson,title={Definice},#1}

% Spolecny box pro zadani a jeho reseni, bez mezery mezi castmi.
% Pouziti:
%   \begin{prikladReseniBox}
%     Zadani ulohy.
%     \tcblower
%     \begin{align*} ... \end{align*}
%   \end{prikladReseniBox}
% Spodni cast ma jemny sedy podklad a delici linku; lze v ni sazet
% vice radku matematiky. Obe casti maji jen jeden spolecny ramecek.
\tcbset{solutionlesson/.style={
  commonlesson,
  enhanced jigsaw,
  breakable,
  colbacklower=black!3,
  segmentation style={solid,draw=black!25,line width=0.35pt},
  fontlower=\small\color{notegray},
  before lower={\noindent\textbf{Řešení:}\enspace},
  leftlower=3mm,rightlower=3mm
}}
\newtcolorbox{prikladReseniBox}[1][]{solutionlesson,title={Příklad},#1}
\newtcolorbox{zadaniReseniBox}[1][]{solutionlesson,title={Zadání},#1}

% Pomucky pro autora pripravy: sedive poznamky; \reseni pro starsi pripravy
\newcommand{\poznamka}[1]{\par\begingroup\small\color{notegray}\itshape
  \textnormal{\textbf{Poznámka pro učitele.}}\enspace #1\par\endgroup}
\newcommand{\reseni}[1]{\par\begingroup\small\color{notegray}
  \textbf{Očekávané řešení:}\enspace #1\par\endgroup}
\newcommand{\blok}[1]{\needspace{5\baselineskip}\subsection*{#1}}
\newcommand{\vyrok}[1]{\uv{#1}}

\nastavPripravu{1. ročník}{Rovnice a nerovnice}{76}{Rovnice s parametrem}
\begin{document}
\blok{Jedna rovnice pro mnoho různých zadání}
V rovnici $(m-2)x=4$ je $x$ neznámá a $m$ označuje zvolené číslo. Pro každé konkrétní $m$ vznikne jiná rovnice. Při $m=0$ je $-2x=4$, při $m=2$ dostaneme $0x=4$. Proto nelze bez podmínky napsat $x=4/(m-2)$.
\begin{definiceBox}
\textbf{Parametr} je číslo, jehož hodnota se v zadání mění mezi jednotlivými případy, zatímco hledáme neznámou $x$. Výsledek zapisujeme zvlášť pro ty hodnoty parametru, při kterých mají použité úpravy odlišné podmínky.
\end{definiceBox}
\begin{definiceBox}
Rovnice $A(m)x=B(m)$ má pro dané $m$ právě jedno řešení $x=B(m)/A(m)$, je-li $A(m)\ne0$; nemá řešení, je-li $A(m)=0$ a $B(m)\ne0$; má všechna reálná čísla za řešení, je-li $A(m)=B(m)=0$.
\end{definiceBox}
\begin{prikladReseniBox}
Vyřešte $(m-2)x=4$ pro všechna $m\in\mathbb R$. Zvlášť ověřte $m=2$.
\tcblower
Pro $m\ne2$ je koeficient u $x$ nenulový, tedy $K_m=\{4/(m-2)\}$. Pro $m=2$ dosadíme do původní rovnice: $0\cdot x=4$, což neplatí pro žádné $x$, tedy $K_2=\varnothing$. Kontrola obecné větve: $(m-2)\cdot4/(m-2)=4$ pouze při $m\ne2$.
\end{prikladReseniBox}
\begin{prikladReseniBox}
Vyřešte $(a-3)x=2a-6$ a porovnejte výjimečný případ s předchozím příkladem.
\tcblower
Pro $a\ne3$ dělíme nenulovým $a-3$ a dostaneme $x=2(a-3)/(a-3)=2$, tedy $K_a=\{2\}$. Při $a=3$ původní rovnice zní $0x=0$ a vyhovuje každé $x$, tedy $K_3=\mathbb R$. Samotný vyloučený jmenovatel proto neznamená „žádné řešení“.
\end{prikladReseniBox}
\poznamka{Nech žáky dosadit konkrétní $a=0$ a $a=3$ do původní rovnice. Rozlišení nulového koeficientu je důležitější než mechanické krácení.}
\blok{Úprava před rozdělením případů}
\begin{prikladReseniBox}
Vyřešte $p(x-1)=2x+3$ s neznámou $x$ a parametrem $p$.
\tcblower
Rozevřeme závorku a shromáždíme členy: $px-p=2x+3$, tedy $(p-2)x=p+3$. Pro $p\ne2$ dostaneme jediné řešení $x=(p+3)/(p-2)$. Pro $p=2$ dosadíme do \emph{původní} rovnice: $2(x-1)=2x+3$, tedy $-2=3$, což je nepravda; $K_2=\varnothing$. Zkouška obecné větve: $(p-2)\cdot(p+3)/(p-2)=p+3$, ekvivalentní upravené rovnici.
\end{prikladReseniBox}
\newpage
\blok{Samostatné procvičení}
\begin{zadaniReseniBox}
\begin{enumerate}[label=\alph*)]
\item Vyřešte $(t+1)x=3$ pro všechna $t$.
\item Vyřešte $k(x+2)=kx+5$ pro všechna $k$.
\item Vyřešte $(s+4)x=2(s+4)$ pro všechna $s$.
\end{enumerate}
\tcblower
\begin{enumerate}[label=\alph*)]
\item Pro $t\ne-1$ je $K_t=\{3/(t+1)\}$; pro $t=-1$ dostaneme $0=3$, tedy $K_{-1}=\varnothing$.
\item Po odečtení $kx$ je $2k=5$. Pro $k=5/2$ platí rovnice pro každé $x$, tedy $K_k=\mathbb R$; pro všechny ostatní $k$ je $K_k=\varnothing$. Nelze „dělit $x$“, protože se vyrušilo.
\item Pro $s\ne-4$ je $K_s=\{2\}$. Pro $s=-4$ se obě strany rovnají nule pro každé $x$, tedy $K_{-4}=\mathbb R$.
\end{enumerate}
\end{zadaniReseniBox}
\blok{Volitelné rozšíření a rezerva}
\begin{zadaniReseniBox}
Vyřešte $(m^2-1)x=m-1$ pro všechna $m\in\mathbb R$ a najděte všechny tři možné typy výsledku.
\tcblower
$m^2-1=(m-1)(m+1)$ je nulové při $m=1$ nebo $m=-1$. Pro $m\notin\{-1,1\}$ je $x=(m-1)/[(m-1)(m+1)]=1/(m+1)$. Při $m=1$ zní původní rovnice $0=0$, tedy $K_1=\mathbb R$. Při $m=-1$ zní $0=-2$, tedy $K_{-1}=\varnothing$. Krácení $m-1$ platilo pouze v obecné větvi.
\end{zadaniReseniBox}
\blok{Domácí práce}
\begin{zadaniReseniBox}
\begin{enumerate}[label=\alph*)]
\item Vyřešte $(b-5)x=b-5$ pro každé $b$.
\item Vyřešte $3x+c=cx+6$ pro každé $c$.
\item Uveďte konkrétní hodnoty parametrů, pro něž příklady a), b) mají jedno, žádné nebo všechna řešení.
\end{enumerate}
\tcblower
\begin{enumerate}[label=\alph*)]
\item Pro $b\ne5$ je $x=1$, tedy $K_b=\{1\}$. Pro $b=5$ vyjde $0=0$, tedy $K_5=\mathbb R$.
\item $(3-c)x=6$. Pro $c\ne3$ je $K_c=\{6/(3-c)\}$. Pro $c=3$ vyjde $0=6$, tedy $K_3=\varnothing$.
\item Např. v a) $b=0$ dává jedno řešení a $b=5$ všechna; žádné tam nenastane. V b) $c=0$ dává jedno řešení a $c=3$ žádné; všechna tam nenastanou.
\end{enumerate}
\end{zadaniReseniBox}
\poznamka{Jádro: motivace a pravidlo 11 min, dvě společné větve 14 min, samostatné úlohy s kontrolou 17 min, organizace 3 min. Třetí řešený příklad lze vybrat podle tempa; rezerva je nad rámec. Domácí práce 15–20 min. Soustavy s parametrem neprobírej.}
\end{document}
